\documentclass{article}
\usepackage[UTF8]{ctex}
\usepackage{amsmath, amssymb, bm, color}
\usepackage{geometry}
\geometry{a4paper, margin=2.5cm}

\title{脉冲神经网络在线学习：关键科学问题与研究推进方案}
\author{}
\date{}

\begin{document}

\maketitle

\begin{abstract}
脉冲神经网络（SNN）凭借其事件驱动、稀疏脉冲通信和生物可解释性，被认为是实现低功耗边缘智能的候选模型。然而，传统的训练算法——尤其是基于时间的反向传播（BPTT）——因需要展开整个时间计算图并存储所有时刻的状态，内存和计算开销随时间步数线性增长，难以在资源受限设备上实现在线学习。近年来，研究者提出了多种在线学习算法，如OTTT、NDOT、S-TLLR和TESS，分别从时序依赖近似、动态核重参数化、STDP启发式资格迹以及本地学习信号生成等角度，实现了与时间步数无关的常数存储复杂度，并在不同程度上降低了计算复杂度。

然而，这些算法仍面临若干根本性挑战：如何在生物合理性与高性能之间取得平衡、如何在存储与计算复杂度之间达到帕累托最优、如何有效应对长期时间依赖、如何与神经形态硬件协同设计，以及如何无缝融合持续学习能力。本文围绕这五个核心问题，结合已有算法的技术特点，提出详细的推进方案，旨在为SNN在线学习的后续研究提供系统性的路线图。
\end{abstract}

\section{引言}

脉冲神经网络（Spiking Neural Networks, SNNs）被誉为第三代神经网络，其神经元通过离散脉冲（spikes）进行通信，在神经形态硬件上具有极高的能效。然而，SNN的监督训练长期受困于脉冲发放函数不可微的问题。基于代理梯度的BPTT（Backpropagation Through Time）虽然通过时间展开实现了有效的梯度传播，但其存储和计算成本均与时间步数 \(T\) 成正比，无法满足在线学习和边缘部署的需求。

为突破这一瓶颈，近年的研究聚焦于设计“时间本地”或“时空本地”的学习规则。OTTT（Online Training Through Time）通过丢弃时序依赖中的重置路径，将时间传播退化为固定指数衰减，实现 \(O(Ln)\) 存储。NDOT（Neuronal Dynamics-based Online Training）进一步用连续动力学推导的动态比值替代固定衰减系数，提升了时序建模精度。S-TLLR（STDP-inspired Temporal Local Learning Rule）引入广义STDP形式的资格迹，融合因果与非因果信息，并依靠反向传播（BP）或直接反馈对齐（DFA）产生学习信号。TESS（Temporally and Spatially Local Learning Rule）则在S-TLLR的基础上，通过本地学习信号生成（LSG）彻底切断了层间误差传播，将时间复杂度从 \(O(Ln^2)\) 降至 \(O(LCn)\)。

尽管这些算法取得了显著进展，但从理论到实际部署仍存在多个尚未解决的关键问题。本文从算法设计、复杂度理论、长期依赖建模、硬件协同设计、持续学习融合等维度，提炼出五个核心科学问题，并逐一给出详尽的推进方案，以期为该领域的后续研究提供参考。

\section{问题一：如何设计同时具备“生物合理性”与“高性能”的本地学习规则？}

\subsection{问题背景与现状}

生物神经系统依赖局部突触可塑性（如STDP、资格迹）进行学习，无需全局误差反向传播。然而，纯生物启发的方法（如无监督STDP）在复杂任务上性能远不及梯度方法。当前在线学习算法在“生物合理性”与“任务性能”之间存在明显断层：

\begin{itemize}
    \item 以OTTT、NDOT为代表的算法，本质上是BPTT的近似，虽性能优越，但其学习规则（尤其是代理梯度的使用）缺乏明确的生物学对应。
    \item 以S-TLLR、TESS为代表的算法，引入了STDP形式的因果/非因果项，更具生物可解释性，但其学习信号的生成仍依赖BP（S-TLLR）或固定随机投影（TESS），与真实的生物机制尚有距离。
\end{itemize}

此外，S-TLLR中非因果项系数 \(\alpha_{post}\) 的最优取值（-1、0、+1）在不同任务上表现迥异，其背后的任务依赖性尚未被系统理解。TESS中固定随机矩阵 \(\mathbf{B}\) 的设计虽然硬件友好，但与BPTT相比仍有约1.4\%的性能差距（以CIFAR10-DVS为例），如何在不牺牲本地性的前提下缩小这一差距，是一个亟待解决的问题。

\subsection{核心研究问题}

\begin{enumerate}
    \item \textbf{Q1.1}：能否建立一个统一的数学框架，将“梯度近似型”（OTTT/NDOT）和“STDP启发型”（S-TLLR/TESS）算法纳入同一形式，并揭示它们之间的等价或互补关系？
    \item \textbf{Q1.2}：S-TLLR/TESS中的非因果项系数 \(\alpha_{post}\) 的最优取值是否与任务的时间结构（如事件密度、脉冲稀疏度）存在确定性关联？能否设计自适应机制，让网络动态调整 \(\alpha_{post}\)？
    \item \textbf{Q1.3}：TESS中LSG所用的固定随机矩阵 \(\mathbf{B}\) 能否被优化（如使用可学习的低秩矩阵、数据驱动的投影或任务相关的正交基），在保持本地性的前提下进一步逼近BPTT的性能？
\end{enumerate}

\subsection{推进方案}

\paragraph{（1）统一框架构建}
从三因子学习规则的一般形式 \(\Delta w_{ij} = \sum_t \delta_i[t] e_{ij}[t]\) 出发，分别写出各类算法的 \(\delta\) 与 \(e\) 的显式表达式。分析OTTT/NDOT中的空间梯度 \(\mathbf{g}_{\mathbf{u}}\) 与S-TLLR/TESS中的本地学习信号 \(\mathbf{v}_m\) 之间的数学关联，以及踪迹 \(\hat{\mathbf{a}}\)、\(\text{tr}\)、\(\mathbf{q}\)、\(\mathbf{h}\) 之间的转换关系。通过变量代换，证明在特定条件下（如忽略非因果项、令 \(\lambda\) 与 \(\gamma\) 相等），S-TLLR可退化为OTTT；而TESS的LSG可视为对S-TLLR中DFA的一种推广。该统一框架将为后续算法设计提供理论指导。

\paragraph{（2）\(\alpha_{post}\) 的自适应机制}
设计可学习的参数 \(\alpha_{post}^{(l)}\)（可逐层或逐任务自适应），将其与网络输出误差相关联，让网络在训练过程中自主决定因果/非因果项的权重。例如，可将 \(\alpha_{post}\) 视为一个可训练标量，通过梯度下降（利用代理梯度）调整，或通过元学习（如MAML）在不同任务上快速适应。同时，设计实验，系统测量不同任务的事件密度、平均发放率、时间相关性等指标，建立 \(\alpha_{post}\) 最优值与这些指标之间的经验映射关系。

\paragraph{（3）\(\mathbf{B}\) 矩阵的优化}
保留LSG的本地性，但将固定的随机 \(\mathbf{B}\) 替换为任务相关的投影矩阵。考虑以下方案：
\begin{itemize}
    \item 使用无监督预训练（如PCA）从数据中提取主成分，作为 \(\mathbf{B}\) 的初始值，并在本地学习过程中保持固定，从而在不引入额外反向传播的前提下获得更具判别性的投影。
    \item 设计可学习的低秩矩阵 \(\mathbf{B} = \mathbf{U} \mathbf{V}^\top\)，其中 \(\mathbf{U}\) 和 \(\mathbf{V}\) 在层内通过本地更新规则调整，但保持层间独立。
    \item 借鉴神经科学中的“随机突触后投影”理论，使用具有准正交性的确定性序列（如Hadamard矩阵或方波函数）替代随机矩阵，以提升投影的稳定性与表征能力。
\end{itemize}

\section{问题二：如何实现“存储复杂度最优”与“计算精度最优”的帕累托平衡？}

\subsection{问题背景与现状}

如前文所述，现有在线学习算法均实现了 \(O(Ln)\) 的存储复杂度，与时间步数 \(T\) 完全解耦。然而，时间复杂度差异显著：

\begin{center}
\begin{tabular}{|c|c|c|c|}
\hline
算法 & 存储复杂度 & 时间复杂度 & 关键操作 \\
\hline
OTTT & \(O(Ln)\) & \(O(Ln^2)\) & 外积 \(\mathbf{g} \cdot \hat{\mathbf{a}}^\top\) + BP \\
NDOT & \(O(Ln)\) & \(O(Ln^2)\) & 外积 + BP（动态核仅需逐元素计算） \\
S-TLLR & \(O(Ln)\) & \(O(Ln^2)\) & 外积 + BP/DFA（\(n^2\)） \\
TESS & \(O(Ln)\) & \(O(LCn)\) & 外积 + LSG（\(Cn\)） \\
\hline
\end{tabular}
\end{center}

除TESS外，其余算法的时间复杂度均为 \(O(Ln^2)\)，因为权重更新本身需要计算 \(n \times n\) 的外积（这是突触数量决定的“物理极限”）。但TESS通过LSG将空间误差传播从 \(O(n^2)\) 降至 \(O(Cn)\)，从而整体上大幅降低了计算量。然而，这一改进是否已触及复杂度下界？是否存在理论上更优的方案？另外，对于不同的任务类型（静态图像 vs. 动态事件流），最优的复杂度-精度折中点可能不同，这需要系统性研究。

\subsection{核心研究问题}

\begin{enumerate}
    \item \textbf{Q2.1}：在SNN在线学习中，是否存在一个理论上的“复杂度下界”——即任何算法都必须至少达到的存储和计算复杂度？特别地，\(n \times n\) 的外积是否不可避免？
    \item \textbf{Q2.2}：能否设计一种算法，在保持 \(O(Ln)\) 存储的同时，将时间复杂度也降至 \(O(Ln)\)（即与 \(n^2\) 完全解耦）？
    \item \textbf{Q2.3}：不同任务类型（静态图像分类、动态手势识别、连续语音流）对复杂度-精度权衡的最优点是否不同？能否设计自适应复杂度调节机制？
\end{enumerate}

\subsection{推进方案}

\paragraph{（1）复杂度下界的信息论分析}
从信息论视角出发，分析突触权重更新所需的最小计算量。考虑一个具有 \(n\) 个输入和 \(n\) 个输出的全连接层，其权重矩阵有 \(n^2\) 个独立参数。在线学习每个时间步都需要更新这些参数，因此任何算法都至少需要 \(O(n^2)\) 次操作来更新所有突触——除非引入结构化稀疏（如卷积的局部连接）。因此，\(O(n^2)\) 可视为“更新所有突触”的固有成本。但我们可以区分“突触更新”和“误差传播”两部分：误差传播可以借助本地信号（如LSG）降为 \(O(Cn)\)，但突触更新本身仍为 \(O(n^2)\)。因此，理论上最优的时间复杂度应为 \(O(n^2 + Cn)\)，当 \(C \ll n\) 时即 \(O(n^2)\)。TESS已将误差传播降至 \(O(Cn)\)，但突触更新部分仍是瓶颈。若采用稀疏连接或结构化权重（如卷积），则可进一步降低。

\paragraph{（2）“极致轻量”算法探索}
借鉴“符号梯度下降”（sign-based gradient descent）或“伪零阶优化”（pseudo-zero-order optimization）的思想，在更新权重时仅使用权重的符号或等级信息，而非完整浮点数梯度。例如，在TESS框架下，可将LSG产生的学习信号 \(\mathbf{v}_m\) 二值化（仅保留正负号），从而将外积计算简化为符号外积，大幅降低计算精度需求，同时可能保持相近性能。此外，可探索“稀疏更新”策略：每时间步仅随机选择一部分突触进行更新，利用随机近似理论保证收敛性。

\paragraph{（3）任务自适应的复杂度调节}
设计一个“复杂度旋钮”——即一个超参数或可配置策略，允许用户根据任务需求（延迟限制、能耗预算、精度要求）动态切换算法内部的计算模式。例如：
\begin{itemize}
    \item 在时间充裕时，使用OTTT/NDOT的完整梯度更新；
    \item 在资源紧张时，切换至TESS的LSG模式；
    \item 在极端受限时，启用二值化或稀疏更新模式。
\end{itemize}
通过在线性能监控，实时调整模式，实现帕累托最优。

\section{问题三：如何有效应对“时空信用分配”的长期依赖问题？}

\subsection{问题背景与现状}

在线学习算法依赖踪迹（eligibility trace）机制来解决时间信用分配问题。踪迹本质上是一个指数衰减的低通滤波器，其记忆长度由衰减系数 \(\lambda\) 控制（如OTTT中的 \(\lambda^{t-\tau}\)）。然而，这种固定指数衰减的“有效记忆范围”约为 \(1/(1-\lambda)\)，当时间步数 \(T\) 远超此范围时，历史信息将指数级消失，导致长期依赖无法被有效捕捉。此问题在S-TLLR和TESS中同样存在，尽管它们通过非因果项部分缓解了时序建模的局限性，但根源未变。

\subsection{核心研究问题}

\begin{enumerate}
    \item \textbf{Q3.1}：现有踪迹机制能否扩展以捕捉更长的时间依赖？是否存在“多时间尺度”踪迹的自然实现？
    \item \textbf{Q3.2}：在极端长时依赖任务（如语音序列、导航决策）中，在线学习算法是否必然逊色于BPTT？差距有多大？
    \item \textbf{Q3.3}：能否借鉴Transformer的注意力机制，设计一种“注意力踪迹”，让网络动态决定关注哪些历史时刻？
\end{enumerate}

\subsection{推进方案}

\paragraph{（1）多尺度踪迹设计}
设计多个不同衰减率的踪迹（\(\lambda_1, \lambda_2, \dots, \lambda_K\)），分别对应短、中、长时间尺度。每个时刻，网络维护一组踪迹：
\[
\hat{a}^{(k)}[t] = \lambda_k \hat{a}^{(k)}[t-1] + s[t], \quad k=1,\dots,K
\]
最终的资格迹为这些多尺度踪迹的加权组合，权重可以通过学习获得（例如，通过一个小的可学习网络或注意力机制）。这样，网络既保留了短时的精确信息，又能积累长时的统计信息，从而应对不同时间尺度的依赖。

\paragraph{（2）注意力式踪迹}
借鉴Transformer中的自注意力机制，设计“踪迹注意力”模块。在每个时间步，计算当前时刻与过去若干个关键时刻的相关性分数，然后加权融合历史信息。具体而言，可维护一个固定长度的“记忆缓冲区”（例如最近 \(M\) 个时刻的脉冲），并利用可学习的查询、键、值向量计算注意力权重，从而动态决定哪些历史时刻对当前更新最重要。虽然这会引入额外的计算开销，但可显著改善长期依赖建模能力。

\paragraph{（3）长时依赖基准实验}
设计一组具有不同时间依赖长度的任务（如合成序列任务、语音命令识别、机器人导航），在这些任务上系统评估OTTT、NDOT、S-TLLR、TESS以及BPTT的性能。测量各算法在不同时间步长下的精度衰减曲线，明确当前在线学习算法的“长时依赖极限”。实验应包括对 \(\lambda\) 超参数的大规模扫描，以揭示衰减系数与任务时间结构之间的匹配规律。

\section{问题四：如何实现“片上学习”的硬件-算法协同设计？}

\subsection{问题背景与现状}

SNN在线学习的终极应用场景是神经形态芯片上的片上学习（on-chip learning）。然而，当前算法设计大多基于通用处理器（GPU）的抽象计算模型，忽略了神经形态硬件的物理约束：低精度定点数运算、有限的片上存储（SRAM）、稀疏脉冲驱动的事件型计算。算法与硬件之间的这种“语义鸿沟”导致即便算法复杂度在理论上很低，实际部署时仍可能面临巨大的效率和精度损失。

\subsection{核心研究问题}

\begin{enumerate}
    \item \textbf{Q4.1}：在硬件约束（如8位定点数、无乘法器、存储受限）下，OTTT/NDOT/S-TLLR/TESS中哪一个最容易实现？哪一个精度损失最小？
    \item \textbf{Q4.2}：能否设计一种“硬件原生”的在线学习算法——其基本操作与神经形态硬件的原生指令集完全对齐（如脉冲累积、移位操作、表查找）？
    \item \textbf{Q4.3}：如何将TESS中的随机矩阵 \(\mathbf{B}\) 在芯片上高效实现，而不消耗大量存储资源？
\end{enumerate}

\subsection{推进方案}

\paragraph{（1）硬件-算法联合仿真平台}
构建一个模拟神经形态芯片硬件约束的仿真框架（如基于PyTorch的量化仿真工具），支持自定义位宽、存储容量、计算精度等参数。在该平台上实现OTTT、NDOT、S-TLLR、TESS，运行统一基准测试，量化不同硬件约束下的精度损失和资源占用。通过该平台，可以识别各算法对硬件最敏感的部分（如S-TLLR的BP信号传播中的矩阵乘法对精度要求高；TESS的LSG则相对鲁棒）。

\paragraph{（2）“脉冲驱动”算法重设计}
重新审视算法中的每一步操作，将其映射为“脉冲驱动”的硬件友好操作：
\begin{itemize}
    \item 踪迹更新 \(\hat{a}[t] = \lambda \hat{a}[t-1] + s[t]\)：在硬件上可通过移位加法和累加实现（\(\lambda\) 可取 \(2^{-k}\)，则乘法变为移位）。
    \item 权重更新 \(\Delta W = \mathbf{v}_m \otimes \mathbf{q}\)：可分解为外积，但若采用二值化或稀疏化，可将乘法简化为逻辑运算。
    \item 学习信号生成（LSG）：\(\mathbf{B} \mathbf{o}\) 可通过固定随机连接（硬件布线）实现，无需存储完整矩阵。
\end{itemize}
重点优化S-TLLR中的BP传播，因其涉及浮点矩阵乘法，可考虑用DFA替代BP以降低硬件开销。

\paragraph{（3）\(\mathbf{B}\) 矩阵的硬件化存储}
利用神经形态芯片中已有的随机连接资源（如Loihi中的突触随机投影功能），将\(\mathbf{B}\)矩阵实现为芯片上的固定随机连接，而非存储在内存中。这样，\(\mathbf{B}\)的存储成本降为 \(O(1)\)，同时投影操作可通过事件驱动的脉冲传播自然完成。若芯片不支持随机连接，则可采用伪随机数发生器在运行时即时生成投影系数，从而避免存储开销。

\section{问题五：在线学习如何与“持续学习”无缝融合？}

\subsection{问题背景与现状}

在线学习通常假设数据来自同一分布（即任务不变），而**持续学习（Continual Learning）** 要求模型在面对新任务时，不遗忘旧任务的知识。SNN的在线学习能力天然适合持续学习场景，因为数据是流式到达的。然而，现有在线学习算法几乎都未考虑任务边界，也无防止灾难性遗忘的机制。当任务发生切换时，权重会被新任务的数据迅速覆盖，导致旧任务性能急剧下降。

\subsection{核心研究问题}

\begin{enumerate}
    \item \textbf{Q5.1}：OTTT/NDOT/S-TLLR/TESS在面对任务切换时，遗忘率如何？哪种算法对连续任务更鲁棒？
    \item \textbf{Q5.2}：能否将现有的持续学习策略（如弹性权重巩固EWC、突触巩固）无缝集成到在线学习的踪迹更新过程中，而无需额外的反向传播？
    \item \textbf{Q5.3}：能否设计一个“睡眠-觉醒”两阶段框架：觉醒时快速在线学习新知识，睡眠时通过重放或生成式模型巩固记忆，同时不额外消耗大量资源？
\end{enumerate}

\subsection{推进方案}

\paragraph{（1）持续学习基准测试}
选择标准持续学习基准数据集（如Split CIFAR-100、Permuted MNIST），在这些数据集上评估四种在线学习算法的遗忘率（即新任务学习后旧任务精度的下降幅度）。同时对比BPTT结合传统持续学习策略（如EWC、LwF）的性能，以建立性能标杆。重点分析各算法在“任务增量”和“类增量”两种场景下的表现差异。

\paragraph{（2）“巩固踪迹”机制}
在现有踪迹更新中引入“巩固项”。借鉴EWC的思想，为每个突触维护一个重要性权重 \(\Omega_{ij}\)，该权重衡量该突触对已学任务的贡献度。在更新资格迹时，加入惩罚项，使得重要性高的突触在后续学习中变化更慢。具体地，可将更新公式修改为：
\[
\Delta w_{ij} = \rho \sum_t \delta_i[t] e_{ij}[t] \cdot (1 - \eta \cdot \Omega_{ij})
\]
其中 \(\eta\) 为巩固强度系数，\(\Omega_{ij}\) 可通过任务结束后计算Fisher信息矩阵的迹来估计，或在每次更新中通过梯度平方的移动平均在线估计。这样，在不引入额外反向传播的情况下，实现了对关键突触的保护。

\paragraph{（3）“睡眠-觉醒”两阶段框架}
设计两阶段学习策略：
\begin{itemize}
    \item \textbf{觉醒阶段}：使用TESS（因其实时性好）对当前数据流进行快速本地学习，实时更新权重。
    \item \textbf{睡眠阶段}：在设备空闲时（如夜间或低负载时段），启动“巩固过程”。这一过程可通过重放存储的少量典型样本（或由生成模型合成的伪样本）来重新激活旧任务，并使用相同的TESS规则进行少量额外训练，但此时学习率降低，同时开启巩固项（如上述 \(\Omega_{ij}\) 的作用）。此外，可探索使用生成式模型（如GNN）在睡眠阶段生成重放样本，以避免存储大量历史数据。
\end{itemize}
该框架将在线学习的“快速性”与持续学习的“稳定性”有机分离，有望在不显著增加在线计算开销的前提下，大幅提升对旧任务的保留能力。

\section{结论}

本文围绕脉冲神经网络在线学习的核心挑战，提炼了五个关键科学问题：生物合理性与高性能的平衡、复杂度-精度的帕累托优化、长期时间依赖建模、硬件-算法协同设计、以及在线学习与持续学习的融合。针对每个问题，我们分析了现有算法（OTTT、NDOT、S-TLLR、TESS）的技术特点与不足，并提出了详细的推进方案。这些方案既包含理论基础（如统一框架、信息论下界），也包含工程实践（如硬件仿真、自适应机制），构成了一套系统性的研究路线图。

未来的研究可以沿着上述方向深入：一方面，从理论上探索在线学习算法的能力边界；另一方面，在实际边缘设备和神经形态芯片上验证算法的有效性。我们相信，随着这些问题的逐步解决，SNN在线学习将真正从学术研究走向实用部署，为低功耗智能系统提供有力的算法支撑。

\end{document}